% source: RKF/theorum/30_stage15_exact_recognition_completion_certificate.md
\hypertarget{stage-1.5-exact-recognition-completion-certificate}{%
\chapter{Stage 1.5 Exact Recognition-Completion Certificate}\label{stage-1.5-exact-recognition-completion-certificate}}

\hypertarget{source-doctrine}{%
\section{1. Source doctrine}\label{source-doctrine}}

This stage consumes the strongest mathematically usable results from the buried Recognition-Seam and morphic-calculus chapters:

\begin{verbatim}
cut and trans-cut morphism;
objects as stabilized morphic lineages;
recovered middle identity;
Recognition-Cauchy completion;
Madhava--Smriti lawful tails;
projection leakage and target faithfulness;
shadow stationarity/convergence versus recognition closure.
\end{verbatim}

The chapters distinguish zero from the pre-cut ground, put the cut inside calculus, and define objects as stabilized recognition classes of morphisms. They also prove by explicit counterexample that Hilbert/norm convergence may coexist with drifting seam memory. The Vedic/Kerala bridge supplies the rule

\[
\text{recognized limit}
=
\text{finite recursion}
+
\text{predeclared correction}
+
\text{Smriti tail}.
\]

The variational chapter adds a useful audit principle: shadow stationarity or conservation does not lift to recognition stationarity unless boundary, phase, memory and holonomy residues also close.

\begin{center}\rule{0.5\linewidth}{0.5pt}\end{center}

\hypertarget{exact-generalized-packet}{%
\section{2. Exact generalized packet}\label{exact-generalized-packet}}

The executable packet is

\begin{verbatim}
proof_lab/clock_free_cut_memory_stage15_examples.py
\end{verbatim}

and the fail-closed tests are

\begin{verbatim}
proof_lab/test_clock_free_cut_memory_stage15.py
\end{verbatim}

All arithmetic is rational.

\begin{center}\rule{0.5\linewidth}{0.5pt}\end{center}

\hypertarget{recovered-middle-identity}{%
\section{3. Recovered middle identity}\label{recovered-middle-identity}}

Two middle states have identical shadows but different cut memories:

\[
B_1=(\operatorname{shadow}B,\phi,\sigma,1),
\qquad
B_2=(\operatorname{shadow}B,\phi,\sigma,0).
\]

Therefore

\[
B_1=_{\rm shadow}B_2,
\qquad
B_1\ne_{\rm RSC}B_2.
\]

Raw composition is rejected. A memory transport declared before observation,

\[
k\mapsto k-1,
\]

restores recovered identity and lawful composition. An identical correction declared after observation is rejected.

This is the exact composition gate needed before finite packets may be compared on one carrier.

\begin{center}\rule{0.5\linewidth}{0.5pt}\end{center}

\hypertarget{shadow-cauchy-is-not-recognition-cauchy}{%
\section{4. Shadow-Cauchy is not Recognition-Cauchy}\label{shadow-cauchy-is-not-recognition-cauchy}}

The sequence

\[
\widetilde x_n=(1,0,1,n)
\]

has constant Hilbert/shadow coordinate, so its shadow distance is identically zero. Its unledgered memory does not form a Recognition-Cauchy sequence.

Under the constitutive memory model

\[
k_n=n
\]

declared before evaluation, the quotient memory is identically zero and the sequence closes in the declared recognition quotient.

Thus forgetting memory and lawfully quotienting memory are mathematically different operations.

\begin{center}\rule{0.5\linewidth}{0.5pt}\end{center}

\hypertarget{madhavasmriti-exact-tail}{%
\section{5. Madhava--Smriti exact tail}\label{madhavasmriti-exact-tail}}

For

\[
r=\frac13,
\qquad
B_\infty=\sum_{j=0}^{\infty}r^j=\frac32,
\]

put

\[
S_n=\sum_{j=0}^{n}r^j,
\qquad
C_n=r^{n+1},
\qquad
\mathsf S_n=\frac{r^{n+2}}{1-r}.
\]

Then, for every tested \(n\),

\[
\boxed{
S_n+C_n+\mathsf S_n=\frac32
}
\]

exactly, and

\[
\mathsf S_{n+1}<\mathsf S_n.
\]

A correction equal to the complete observed residual is rejected because it is defined from the held-out limit after evaluation.

\begin{center}\rule{0.5\linewidth}{0.5pt}\end{center}

\hypertarget{morphic-object-stabilization}{%
\section{6. Morphic object stabilization}\label{morphic-object-stabilization}}

Two different locally closed paths produce the same recovered final state:

\begin{verbatim}
path a: shadow +1, phase +1/6, scale x2, memory +1;
        shadow +1, phase -1/6, scale x1/2, memory -1;

path b: shadow +2, phase 0, scale x1, memory 0.
\end{verbatim}

Both stabilize at

\[
(2,0,1,0).
\]

A negative-control path returns the same shadow but leaves memory one. It is not the same morphic object in RSC.

\begin{center}\rule{0.5\linewidth}{0.5pt}\end{center}

\hypertarget{recognition-complete-five-memory-limit}{%
\section{7. Recognition-complete five-memory limit}\label{recognition-complete-five-memory-limit}}

The recognized channel is fixed and has reference floor

\[
R=I_5.
\]

The target-visible memory amplitudes converge to

\[
\left(\frac12,\frac23,\frac34,\frac45,\frac9{10}\right),
\]

while a sixth blind memory coordinate has amplitude

\[
\varepsilon_n=\frac1{n+2}\longrightarrow0.
\]

Thus the finite packets have a genuine hidden complement, but the limit does not.

The target-faithfulness residual contracts from

\[
\frac1{10}
\quad\text{at }n=8
\]

to

\[
\frac1{18}
\quad\text{at }m=16.
\]

The limit memory rank is exactly five. The finite target seam matrix at \(n=8\) has top

\[
\lambda_{\max}(B_8)=\frac{654481}{810000},
\]

and the exact outward transfer error is

\[
e_8=\frac{1619}{810000}.
\]

Therefore

\[
\boxed{
\lambda_{\max}(B_8)+e_8
=
\frac{81}{100}<1.
}
\]

The limit relative gap is exactly

\[
\boxed{
1-\lambda_{\max}(B)=\frac{19}{100}.
}
\]

This is a generalized exact finite-to-infinite certificate. It is not an RH certificate.

\begin{center}\rule{0.5\linewidth}{0.5pt}\end{center}

\hypertarget{expected-status}{%
\section{8. Expected status}\label{expected-status}}

\begin{verbatim}
PASS_CLOCK_FREE_CUT_MEMORY_STAGE15
\end{verbatim}

with six fail-closed unit tests and a byte-stable JSON certificate.

\begin{center}\rule{0.5\linewidth}{0.5pt}\end{center}

\hypertarget{gold-audit-completed}{%
\section{9. Gold audit completed}\label{gold-audit-completed}}

All nineteen uploaded calculus chapters and the Vedic/Kerala bridge note have now been read. Their useful material has been sorted into three classes.

\hypertarget{adopted-foundation}{%
\subsection{Adopted foundation}\label{adopted-foundation}}

\begin{verbatim}
cut/trans-cut calculus;
clock-free transition differential;
composition residues;
recovered identity;
Recognition-Cauchy completion;
Madhava--Smriti tails;
projection leakage;
spectral/Fredholm shadow discipline;
recognition variational and Noether residue audit.
\end{verbatim}

\hypertarget{representation-charts}{%
\subsection{Representation charts}\label{representation-charts}}

\begin{verbatim}
UGD phase-scale-memory coordinates;
circular/hyperbolic R-K chart;
exponential change families;
PDE recognition-regime charts;
Crown connection notation.
\end{verbatim}

\hypertarget{not-promoted-without-further-representation}{%
\subsection{Not promoted without further representation}\label{not-promoted-without-further-representation}}

\begin{verbatim}
undeclared morphic commutator constants;
formal Crown components lacking operators/domains;
curvature or physical interpretations inferred from notation alone;
recognition-index tuples without a declared target group.
\end{verbatim}

\begin{center}\rule{0.5\linewidth}{0.5pt}\end{center}

\hypertarget{next-stage-after-execution}{%
\section{10. Next stage after execution}\label{next-stage-after-execution}}

After the user confirms the six tests and byte-stable JSON comparison, the next generalized theorem will be the Cut-Variational Minimal Observer Theorem.

Its purpose is to derive target faithfulness from the cut-memory action

\[
\mathcal A_\Sigma(P)
=
\|(I-P)M\|^2,
\]

rather than assuming a five-label observer. It will prove:

\begin{verbatim}
zero action iff the observer contains the complete adverse memory range;
minimum faithful observer rank equals the target-relevant cut-memory rank;
source-restriction blind dimension five implies five labels are necessary and sufficient;
recognition-complete convergence transfers finite variational faithfulness to the limit.
\end{verbatim}

Only after that universal theorem passes will the completed-Weil event family be introduced.

\begin{center}\rule{0.5\linewidth}{0.5pt}\end{center}

\hypertarget{claim-boundary}{%
\section{11. Claim boundary}\label{claim-boundary}}

\begin{verbatim}
Stage-1.5 exact examples                         WRITTEN
Stage-1.5 local dry run                          PASSED BY DEVELOPMENT AUDIT
user machine tests                               PENDING
byte-stable user certificate                     PENDING
Cut-Variational Minimal Observer Theorem         NEXT AFTER USER PASS
completed-Weil event lift                        NOT YET STARTED
RH                                                NOT CLAIMED
\end{verbatim}

\begin{flushright}\small\texttt{RKF/theorum/30\_stage15\_exact\_recognition\_completion\_certificate.md}\end{flushright}
